% =====================================================================
%  Thème 5 — Angle et identités remarquables — MOYEN — Corrigé
% =====================================================================
\documentclass[11pt,a4paper]{article}
\usepackage[utf8]{inputenc}
\usepackage[T1]{fontenc}
\usepackage{lmodern}
\usepackage[french]{babel}
\usepackage{amsmath,amssymb,mathtools}
\usepackage{xcolor}
\usepackage{enumitem}
\usepackage{fancyhdr}
\usepackage[a4paper,margin=2.1cm,headheight=26pt,headsep=10pt]{geometry}

\definecolor{primary}{RGB}{222,70,80}
\definecolor{primarydark}{RGB}{150,30,42}
\definecolor{excol}{RGB}{0,135,90}

\pagestyle{fancy}\fancyhf{}
\fancyhead[L]{\small\textcolor{primarydark}{\textbf{Fiche 2} — Calcul vectoriel}}
\fancyhead[R]{\small\textcolor{primarydark}{\textsc{Thème 5 — Moyen — Corrigé}}}
\fancyfoot[C]{\small\thepage}
\renewcommand{\headrulewidth}{0.8pt}

\newcommand{\vc}[1]{\overrightarrow{#1}}
\providecommand{\norm}[1]{\left\lVert #1\right\rVert}
\newcommand{\cor}[1]{\medskip\noindent{\color{primarydark}\bfseries\large Corrigé #1.}\par\smallskip}
\newcommand{\rep}{\textcolor{excol}{$\blacktriangleright$}\ }

\begin{document}
\begin{center}
{\LARGE\bfseries\color{primarydark}Angle et identités remarquables — Corrigé Moyen}\\[2pt]
{\color{primary}\rule{0.55\linewidth}{1.1pt}}
\end{center}
\vspace{2mm}

\cor{1}
\begin{enumerate}[label=\textbf{\alph*)},leftmargin=2em,itemsep=2pt]
  \item $\vec u\cdot\vec v=1\times4+2\times0=4$ ; $\norm{\vec u}=\sqrt{1+4}=\sqrt5$ ;
        $\norm{\vec v}=\sqrt{16}=4$.
  \item $\cos\theta=\dfrac{4}{\sqrt5\times4}=\dfrac{1}{\sqrt5}=\dfrac{\sqrt5}{5}\approx0{,}447$, donc
        $\theta=\arccos\!\bigl(\tfrac{\sqrt5}{5}\bigr)\approx 63{,}4^\circ$.
\end{enumerate}

\cor{2}
$\vec u\cdot\vec v=0\times2+2\times2\sqrt3=4\sqrt3$ ; $\norm{\vec u}=2$ ;
$\norm{\vec v}=\sqrt{2^2+(2\sqrt3)^2}=\sqrt{4+12}=4$. Donc
\[
  \cos\theta=\frac{4\sqrt3}{2\times4}=\frac{4\sqrt3}{8}=\frac{\sqrt3}{2}\ \Longrightarrow\
  \theta=30^\circ.
\]
\rep Quand le cosinus vaut $\tfrac{\sqrt3}{2}$, l'angle est \emph{exactement} $30^\circ$ : pas besoin
de calculatrice.

\cor{3}
\begin{enumerate}[label=\textbf{\alph*)},leftmargin=2em,itemsep=2pt]
  \item $\vec u+\vec v\,(3;4)$, $\norm{\vec u+\vec v}=\sqrt{9+16}=\sqrt{25}=5$.
  \item $\norm{\vec u}^2=5$, $\norm{\vec v}^2=10$, $\vec u\cdot\vec v=2\times1+1\times3=5$, donc
        $\norm{\vec u+\vec v}^2=5+10+2\times5=25$ : on retrouve $5$.
\end{enumerate}

\cor{4}
\begin{enumerate}[label=\textbf{\alph*)},leftmargin=2em,itemsep=2pt]
  \item $37=9+16+2\,\vec u\cdot\vec v \Rightarrow 2\,\vec u\cdot\vec v=37-25=12 \Rightarrow
        \vec u\cdot\vec v=6$.
  \item $\cos(\vec u,\vec v)=\dfrac{\vec u\cdot\vec v}{\norm{\vec u}\,\norm{\vec v}}
        =\dfrac{6}{3\times4}=\dfrac12$, donc l'angle vaut $60^\circ$.
\end{enumerate}
\rep L'identité remarquable permet de trouver $\vec u\cdot\vec v$ \emph{sans connaître les composantes},
à partir des seules normes.

\cor{5}
\begin{enumerate}[label=\textbf{\alph*)},leftmargin=2em,itemsep=2pt]
  \item $\vc{AB}\,(3;1)$, $\vc{AC}\,(1;4)$ ; $\vc{AB}\cdot\vc{AC}=3\times1+1\times4=7$ ;
        $\norm{\vc{AB}}=\sqrt{10}$, $\norm{\vc{AC}}=\sqrt{17}$.
  \item $\cos\widehat{BAC}=\dfrac{7}{\sqrt{10}\,\sqrt{17}}=\dfrac{7}{\sqrt{170}}\approx0{,}537$, donc
        $\widehat{BAC}\approx 57{,}5^\circ$.
\end{enumerate}
\rep Un angle géométrique $\widehat{BAC}$ est l'angle entre les vecteurs $\vc{AB}$ et $\vc{AC}$ (tous
deux issus de $A$).

\end{document}
